\scanpage{153}
\mspage{476}
% position: upper note 108 1
\(({}^{108}1)\) (6 octobre) Je veux dire que les forces de nature répressive qui ont joué
dans ma vie, semblent prendre surtout, sinon exclusivement, l'une de ces deux
formes spécifiques: enterrement du passé, et mise en avant de mes traits ``\unsure{virils|virils''}
au détriment de mes traits ``féminins''. Je n'entends nullement dire que ces deux
``forces'', de nature répressive l'une et l'autre (c'est à dire, visant à un ``re-\linebreak
foulement'', à un escamotage d'une certaine réalité), soient les seules qui aient
``dominé ma vie''! Ce serait oublier tout l'aspect non égotique de mon être, la
pulsion de connaissance s'exprimant aussi bien au niveau de corps que de \unsure{l'espri|l'esprit}
(Voir à ce sujet notamment ``Mes passions'', section n° 35.)

% position: middle paragraph
Même parmi les forces stucturant le moi, émanation du ``patron'' donc, il en
est au moins une, de nature non répressive par elle-même, bien antérieure aux
forces de répression et dont le rôle dans ma vie a été\quad plus essentiel encore:
c'est l'identification à mon père, qui a été comme ``le coeur paisible et \unsure{puissan|puissant''}
du sentiment de ma propre force. Cette identification n'allait nullement dans le
sens de l'exaltation de certaines valeurs ou qualités (viriles disons) au détri-\linebreak
ment d'autres (``féminines''). Indépendamment des valeurs professées par mon père,
sa personne (jusqu'en 1933, où un basculement a eu lieu en lui (*)), a été em-\linebreak
preinte d'un fort équilibre yin-yang, où l'intuition et la spontanéité n'avaient
pas une moindre part que l'intellect et la volonté.

% position: lower paragraph
Enfin, comme autre ``force'' importante de nature égotique, intimement liée,
elle, aux mécanismes répressifs (ou pour mieux dire, de nature ``répressive''
elle-même), il convient encore de compter la sempiternelle \mbox{v a n i t é}, dont
le rôle a été aussi lourd dans ma vie que dans celle de quiconque. Mais cette
``force''-là est de nature si universelle, tout comme le rôle dominant qu'elle
joue dans la vie de chacun (sous une forme plus ou moins grossière ou subtile),
qu'il n'y a guère lieu de l'inclure expressément, dans un relevé des formes spé-\linebreak
cifiques que prennent chez une personne les forces et mécanismes qui structurent
le moi, et donnent à celui-ci sa physionomie particulière et son assise.

% position: footnote rule
\figure{\begin{tikzpicture}\draw (0,0)--(3,0);\end{tikzpicture}}
% position: footnote
(*) Chose remarquable\inserted{,}ce ``basculement'' chez mon père (agé alors de 43 ans) s'est
fait vers un état \mbox{s u p e r - y i n}, vers une sorte de passivité de pacha, en
étroite connivence avec ma mère, jouant un rôle super-yang. Celle-ci l'a pris en
charge\quad lieu et place de ses enfants.(Ils sont largués aux ``profits et per-\linebreak
tes'', jusqu'en 1939 tout au moins, l'année où sous la pression des événements et
à son corps défendant, elle finit par me reprendre auprès d'elle...) Cette rela-\linebreak
tion de dépendance de mon père et de renversement des rôles yin-yang entre mes
parents, a duré jusqu'à la disparition de mon père en 1942.
% position: left foot short mark
\margin{\unsure{-}}
% position: foot pencil
\margin{152}
